\section{Menu fees and the two response budgets}
\label{sec:ex-lattice-rg-multisource}

This section expands the two quantitative suppliers consumed by the selected
one-step construction. Work at one selected adjacent pair of levels and
retain its schedule, menu family, tuned representation, and activity model.
All sums below use the selected countable menu or field index sets; no new
selection of menus or rates is made.

\subsection{The paid fragrance response}
\label{subsec:ex-lattice-rg-usable}

Let $i$ be a fragrance label. Its order is $n_i$ and its spatial labels are
$u_{i,1},\ldots,u_{i,n_i}$. Write $c_i$ for the corresponding Ursell
coefficient and $t_{i,\ell}$ for the selected tuned component, with
$\ell$ running over the distinguished base label and the $n_i$ spatial
labels. The fee majorant is
\begin{equation}
 \ell_i=\frac{|c_i|}{(n_i+1)!}
          \prod_{\ell\in\{0,1,\ldots,n_i\}}t_{i,\ell}.
 \label{eq:ex-lattice-multi-lace}
\end{equation}
A switching observable is paid at cost $c$ when its coefficient series is
controlled by a nonnegative majorant with total sum at most $c$, uniformly
for all allowed interpolation parameters and configurations. The trained
fragrance response is
$s_i(U,t,F)=Y_{i}(\operatorname{color}(U,t,F))$; the paid observable is
\[
 \tau_i(U,t,F)=s_i(U,t,F)-s_i(U,0,\varnothing),
\]
where $F$ is a finite set of selected menu indices. The product formula for
$s_i$ writes it as its Ursell coefficient divided by $(n_i+1)!$ times the
product of the component votes. The product has $n_i+1$ factors: one distinguished base position and the
$n_i$ spatial positions. The switching payment for a
product is as follows. For a finite marker set $I$, a score family $g_t(A)$
represents
\[
 f_t(F)=\sum_{A\subseteq F}t^{|A|}g_t(A),\qquad 0\leq t\leq1.
\]
A payment certificate supplies a nonnegative majorant $b(A)$, independent
of $t$ and the configuration, with $|g_t(A)|\leq b(A)$,
$\sum_A b(A)\leq c$, and convergence of each score as $t\downarrow0$.
If $g_t$ and $h_t$ represent two paid observables, the score for their
product is
\[
 (g_t\star_t h_t)(C)=\sum_{A\cup B=C}t^{|A\cap B|}g_t(A)h_t(B).
\]
Expanding the product and regrouping by $C=A\cup B$ gives the original
weight, since $|A|+|B|=|A\cup B|+|A\cap B|$. The coefficient bound at
$t=1$ is the same convolution of the two nonnegative majorants; its total
sum is the product of their sums. There are finitely many scores, so
continuity at $t=0$ is preserved. Iterating over the $n_i+1$ component votes
and multiplying by the scalar coefficient gives the fee $\ell_i$ in
\eqref{eq:ex-lattice-multi-lace}. The zero-menu term is paid at the same
cost; the triangle inequality for the two score majorants therefore pays
$\tau_i$ at cost $2\ell_i$.

The selected forest estimate gives a ready cluster bound at exponent
$a+e$, with rooted budget $P_0$ and degree ratio $0\leq\rho<1$.
Here $a$ is the selected earned exponent and $e$ the selected allowance;
in the one-step data, $a+e=\operatorname{earn}+w=r_4$. Fix a site $z$ and
mark each underlying component whose full support contains $z$. The
zero-depth cluster-tail estimate sums all connected clusters containing
at least one marked component. Its geometric series is
$\sum_{m\geq0}\rho^m=(1-\rho)^{-1}$, so the total cluster charge at $z$
is at most $P_0/(1-\rho)$. For each such cluster, the connected-coefficient
cover inequality bounds its Ursell coefficient times the support currency
by the fee majorant of its union support. The resulting fragrance fee
ledger is exactly
\begin{equation}
 \sum_{i:z\in R_i}\ell_i
       \exp\!\left(r_4\operatorname{coverNumber}(B,R_i)\right)
 \leq \frac{P_0}{1-\rho},
 \label{eq:ex-lattice-multi-lucent}
\end{equation}
where $R_i$ is the full fragrance support and $B$ is the selected lower
radius. Here $P_0$ is the pledged first-row allotment, and $\rho$ is the
upward sail loss. This is the zero-depth specialization of the geometric
cluster-tail estimate. No extra exponent margin is used.
The coefficient sequence is summable with an explicit finite-volume factor.
The selected conditions give $r_4\geq0$. For nonnegative fees on nonempty
supports, sum the rooted bound over all sites in the finite lower torus.
Each fee contributes at at least one site, with its exponential factor at
least one. There are $|\Lambda_c|$ sites, so
\[
 \sum_i\ell_i\leq |\Lambda_c|\frac{P_0}{1-\rho}.
\]
This is the source's nonempty-support counting argument. Its finite-volume
factor is needed for summability; it is not absorbed into the rooted budget.
The rooted bound prices supports, while this summability justifies the
countable expansion.

The subtraction cost is $2\ell_i$. Multiplying the rooted fee sequence
by two scales its rooted sum and budget by two. The selected stage identities identify the resulting exponent and
budget as
\[
 r_4=e_{\mathrm{earn}}+w,\qquad
 d_4=P_0(1-\rho)^{-1},
\]
so the response fee ledger has exponent $r_4$ and budget $2d_4$. Its zero
coefficient at the empty marker mode follows from the ways/color identity:
when the marker set is empty, every way has the same value as at interpolation
parameter zero, so the trained response equals its subtracted zero-menu
value. Nonemptiness of each menu mode is used separately in the nilpotence
calculation below. The exact locality statement is a noninterference implication. For each
menu label $i$, marker site $u$ in its crest image, and activity label $j$,
if $u$ is outside the support of the $j$th forced activity term, erasing
switch $i$ leaves that activity term unchanged. Consequently a term can
depend on switch $i$ only when its support contains the entire selected
crest image. The implication passes to the absolutely convergent sum, with
no cancellation argument.

For the combined observable, the selected cash identity first holds as an
identity of switched scalar values, for every $t\in[0,1]$, finite marker
set $F$, and configuration $U$:
\begin{equation}
 e^{E(U)}e^{\sum_i\tau_i(U,t,F)}
 \operatorname{gas}_{\mathcal A}\bigl(K(t,F)\bigr)(U)
 =\operatorname{push}(\operatorname{buy}(t,F))(U).
 \label{eq:ex-lattice-multi-closed}
\end{equation}
Here $K(t,F)$ is the final response code, and $\mathcal A$ is the same
activity family before and after the push. To verify the identity, let
$L_0,L_1$ be the logarithm sums of the entrance and branched responses and
let $G_0,G_1$ be their gas polynomials. The switching calculation gives
$G_1=G_0e^{L_0}$ and
$\operatorname{gas}_{\mathcal A}(K)=G_1e^{L_1}$. The base purchase identity
is $\operatorname{push}(\operatorname{buy})=C(U)M(U)G_0(U)$, where $C$
is the selected overhead and $M$ the homestead multiplier. The cash
assembly gives $M=e^{\sum_i s_i(U,t,F)}$ and the zero-menu multiplier is
$e^{t_0}=\mathsf P_{\rm vac}(U)$. Finally, the selected exponent identity is
\[
 e^{E(U)}e^{L_0+L_1}=C(U)\mathsf P_{\rm vac}(U).
\]
Multiplying the switching identities and substituting this exponent identity
reduces the left side of \eqref{eq:ex-lattice-multi-closed} to
$C(U)M(U)G_0(U)$, which is the base purchase. The paid fees are summable,
so the logarithm series converges absolutely and each switched scalar product
is well-defined. For $t>0$, the triangular evaluation map
$g\mapsto(f(F)=\sum_{A\subseteq F}t^{|A|}g(A))$ is injective: induction
on $|F|$ recovers $g(F)$ after the proper-subset terms are known. Thus the
switched scalar identity determines equal score families for every $t>0$;
letting $t\downarrow0$ under the uniform summable score bounds gives equal
limiting scores. To obtain the nilpotent polynomial identity, for a family
of scores $g_t(A)$ and square-zero elements $z_i$, take
$\sum_A(\lim_{t\downarrow0}g_t(A))\prod_{i\in A}z_i$. Addition is
preserved by linearity. Multiplication is preserved because the score
convolution has weight $t^{|A\cap B|}$, whose limit is zero whenever the
same marker occurs in both factors; the corresponding product then vanishes
by $z_i^2=0$. The exponential follows by applying this product rule to its
power series. For the countable logarithm sum, the uniform summable fee
majorants permit dominated convergence coefficient by coefficient and then
summation over markers. This establishes the coefficientwise closed-response
identity used in the first paid response certificate; scalar equality alone
would not suffice.

The full response certificate also includes a paid gas response and a
locality statement. These are distinct obligations: the rooted fragrance
ledger above does not imply the gas budget. The source proves the first by
combining the selected hard and soft budgets into a rate price, then
harvesting it through the fused response. The locality statement follows
from the explicit crest support of each menu contribution.

\subsection{Pricing the second response row}
\label{subsec:ex-lattice-rg-crowd}

At level $m$, let $P$ be the lower budget point, $\Xi$ the selected
parameter profile, $T_0$ the shelved activity, $T_1$ the branched
activity, and $\alpha$ the selected kinetic parameter. Write $S,\sigma,C$
for the clip, width and path count in $\Xi$, and
$H=\operatorname{credit}_G(\Xi)\delta_E(\beta)$. The time profile and
budget point satisfy the selected matching relation. The ordinary field
request is
\begin{equation}
 f=\operatorname{softNeed}_G(P,S,\sigma,C,\alpha,H)
 \label{eq:ex-lattice-multi-softneed}
\end{equation}
By definition this is
$f=\operatorname{softReset}_G(P,H)+\operatorname{hitNeed}(P,\alpha,H)
\operatorname{lowHard}(P,S,\sigma,C)$.
Write $R=\operatorname{raiser}(P)$ for the raised point. It changes only
the support exponent, so its radius agrees with the base radius:
$B_+=B=P.B$. Retain the
same selected profile $\Xi$ for its clip, width, path count, and credit.
The raised final demand is the base-point price which pays a vapor receipt
at $R$:
\begin{equation}
 f_+=F(H)+N(P,\alpha,H)Q(P,S,\sigma,C),
 \label{eq:ex-lattice-multi-raised-demand}
\end{equation}
Let $N(P,\alpha,H)=\operatorname{hitNeed}(P,\alpha,H)$ and define the
selected reset and soft-column prices by
\begin{align*}
 F(H)&=\mathcal P(1,4b+2,0,1/2,\operatorname{resetLoss}(G,P),
                          H,e+2,81),\\
 Q(P,S,\sigma,C)&=\mathcal P(1,b,0,\sigma/8,1,
     d_{\mathrm{sig}}L_{\mathrm{low}}(S,C),e+2,C_{\mathrm{low}}),
\end{align*}
where $b=P^B$, $e=P.\mathrm{supportExponent}$,
$d_{\mathrm{sig}}$ is the number of signed directions, and $C_{\mathrm{low}}$
is the fixed low-column count. The footprint price used here is the explicit
function
\[
 \mathcal P(P,W,k,\epsilon,C,\gamma,a,v)
 =\frac{e^{av}\gamma C\,P\,\operatorname{ballCost}(W)
         \,\epsilon^{-(4+k)}(4+k)!e^{\epsilon}}
        {1-e^{-\epsilon}}.
\]
Thus $f_+=F(H)+N(P,\alpha,H)Q(P,S,\sigma,C)$ is the reset plus
hit-weighted soft-column cost. These are the two terms in the enlarged-
support vapor estimate at the raised point $R$; in general $f_+\ne f$.
The estimate is used with the profile's clip, width and path count, and its
credit $H=\operatorname{credit}_G(\Xi)\delta_E(\beta)$. The matching
relation supplies equality of power, order and support exponent between
$P$ and $\Xi$. The estimate requires the selected local table, bounds, width and
decent-row conditions, the room condition, and a coarse torus side length
greater than $20$. In the source
both footprint prices use the base budget's support exponent plus two;
raising the budget changes the
target receipt, not the argument at which the numerical price is evaluated.
There are two distinct applications of the arrival certificate in the
selected proof. The precursor is the purchased entrance response, of cost
$z_0=2P_Q\exp(-E_0)$, where $P_Q$ is the selected pledge and $E_0$ its
marking exponent. The gain inequality gives $2P_Q\leq1$, so
$0\leq z_0\leq\exp(-E_0)$. For the ordinary path, the column-tail
receipt gives $0\leq p\leq\beta_F^{-1/8}$; the first arrival clause
prices the precursor and the finite-product harvesting estimate transfers
it to the branched activity at $\operatorname{stems}$. For the raised path, the first-row column receipt is instead
$0\leq p_3\leq\exp(\chi_0)\beta_F^{-1/8}$, with $\chi_0$ denoting the selected cover ticket. This is the soft allowance paired with its raised precursor. The
final vapor estimate is a separate local estimate and has a different
budget, stated below.

For the final row, the second arrival clause is used with the exact same
budget, clip, width, path count, and schedule identity. Its hypotheses are
$|\alpha|\leq\alpha_{\max}$, equality of the current parameters with the
selected profile, and $f\geq0$. It returns a rate price $\kappa$ for the
pair of goals
\[
 \operatorname{straw}=2\,\operatorname{lauge}(L,\mu/4,\beta,1/4),
 \qquad
 \operatorname{stems}=2\,\operatorname{lauge}(L,\mu/4,\beta,1/2).
\]
By the definition of the exponential taper,
\[
 \operatorname{stems}=2\,\operatorname{lauge}(L,\mu/8,\beta,1),\qquad
 \operatorname{straw}=2\,\operatorname{lauge}(L,\mu/8,\beta,1/2).
\]
This exact factor of two is kept in the payment: no volume factor is
removed from the row budget.

The rate certificate has hard goal $\operatorname{stems}$, soft goal $f$,
and target toll $\operatorname{straw}$. Its two numerical conditions are
\begin{align}
 &(2B+1)^4\bigl(e^\kappa\operatorname{stems}+2f\bigr)\leq1,
 \label{eq:ex-lattice-multi-capital}\\
 &\operatorname{Quote}(\kappa,
      e^\kappa\operatorname{stems}+2f)\leq\operatorname{straw}.
 \label{eq:ex-lattice-multi-rate}
\end{align}
Here $B$ is the integer budget radius. The rate also requires
$\kappa\geq0$, $0\leq\operatorname{stems}$, and $0\leq f\leq1$.
The first inequality pays the fourth-degree local inflation and the soft
reserve; the second is the selected financial price. The arrival certificate
supplies all these fields. To pass a paid precursor to the widened activity family, use the selected
injection $u$ from old activity labels to the new drawer labels. Put the old
code at label $u(i)$ and set it to zero at every label outside the image.
The widening identity says the forced field at $u(i)$ is exactly the old
forced field at $i$, and all forced fields outside the image vanish. Thus
the score majorant is extended by zero with the same total price; the
physical supports at retained labels are unchanged, so the rooted hit
receipt is transported without any site map. The payment's support-exponent
allowance is the one supplied by the widening step.

Set $b=\operatorname{stems}$ and $w=f$. The rate inflation is
$\Gamma=e^\kappa b+2w$. By definition,
$\operatorname{Quote}(\kappa,\Gamma)=e^{-\kappa}(e^\Gamma-1)$.
The capital and quote inequalities above therefore give exactly the
hypotheses and site output of the hard/soft fusion lemma in
Section~\ref{sec:soft-log-fusion}:
\[
 (2B+1)^4\Gamma\leq1,\qquad
 e^{-\kappa}(e^\Gamma-1)\leq\operatorname{straw}.
\]
That section proves the gas identity
\eqref{eq:fusion-gas-identity} by a weight-preserving bijection of
compatible families. It also proves the priced site estimate
\eqref{eq:fusion-site-output} from $h_B\Gamma\leq1$. In the selected
construction, the widening identity transports the paid hard response, and the
finite-product harvesting estimate applies these results to the hard
precursor at $\operatorname{stems}$ and the soft logarithmic response at
$f$. The rate return inequality is exactly
$e^{-\kappa}(e^\Gamma-1)\leq\operatorname{straw}$, so
\eqref{eq:fusion-site-output} gives the purchased site budget. All overlap
cases are accounted for by the compatible-family bijection, including
eligibility and exact flag coverage.

For the ordinary path, the second arrival clause prices the request
\eqref{eq:ex-lattice-multi-softneed} directly at the straw row, using the
selected profile equality, exact schedule identity, and within-step bound on
$|\alpha|$. For the raised path, the enlarged-support local vapor estimate pays
$f_+$ from \eqref{eq:ex-lattice-multi-raised-demand}; the
cover-weighted allowance $p_3$ belongs to its earlier precursor receipt,
not to this final vapor bound. The final vapor receipt is attached to the
raised budget, while its second rate clause prices at the base budget point
with the selected profile's clip, width, paths and credit. In particular,
with raised rate $\kappa_+$ and the common radius $B_+=B$ it requires
\[
 (2B_++1)^4(e^{\kappa_+}\operatorname{stems}+2f_+)\leq1,
 \qquad e^{-\kappa_+}(e^{e^{\kappa_+}\operatorname{stems}+2f_+}-1)
       \leq\operatorname{straw}.
\]
The proof first widens the support and then harvests the finite product. In
both cases the payment is transported to the result activity by the
finalization identity.
The ordinary paid-response certificate uses the first path; a separate
raised branch of the construction uses the second path. Neither path follows from
the other by changing only the exponent.

\subsection{Constructing the per-box response from the paid rows}
\label{subsec:ex-lattice-rg-canopy}

For the per-box construction, let $z_i$ be the normalized polynomial
supported on the nonempty marker mode $S_i$, with coefficient $a_i/\phi_i$,
where $a_i$ is the menu charge and $\phi_i=\operatorname{cost}_r(S_i)|a_i|$.
Let $stamp_i$ be the polynomial on that mode with coefficient $a_i$. Then
\begin{equation}
 [z_i]_{\varnothing}=0,\qquad z_i^2=0,\qquad
 \operatorname{price}_r(z_i)\leq1,\qquad \phi_i z_i=stamp_i.
 \label{eq:ex-lattice-menu-singleton}
\end{equation}
If $\phi_i=0$, positivity of the cost on a nonempty mode forces $a_i=0$,
and these statements hold with both polynomials zero. If $\phi_i>0$, they
follow directly from the definition of the weighted price and the fact that
the mode is not disjoint from itself.

For each activity label $j$, configuration $U$, and $t\in[0,1]$, write
$g_{j,t}(A,U)$ for the switching score of the selected code. The paid
certificate gives a nonnegative summable majorant $b_j(A)$ with
$|g_{j,t}(A,U)|\leq b_j(A)$, and each score has a limit as $t\downarrow0$.
Define the activity polynomial
\[
 F_j(U)=\sum_{A}\left(\lim_{t\downarrow0}g_{j,t}(A,U)\right)
                       \prod_{i\in A}z_i.
\]
Its weighted price is bounded by $\sum_A b_j(A)$: submultiplicativity of
weighted price gives
\[
 \operatorname{price}_r\Bigl(\prod_{i\in A}z_i\Bigr)\leq1.
\]
The triangle inequality then gives the bound first for finite partial
sums and then for their absolutely convergent limit. A product of two
normalized menu factors either vanishes by a repeated marker or has support
of size at least two. Therefore the coefficient of a singleton marker in
$F_j$ comes only from singleton score sets $A=\{i\}$.

The exact local jet receipt, proved in
Section~\ref{subsec:ex-lattice-menu-jet-proof}, is
\begin{equation}
 \lim_{t\downarrow0}g_{j,t}(\{i\},U)
       =\phi_i D_{i,j}(U),
 \label{eq:ex-lattice-menu-jet-input}
\end{equation}
where $D_{i,j}$ is the derivative of the selected one-menu final response
at activity label $j$ and configuration $U$. Using it in the coefficient formula above and
then using $\phi_i z_i=stamp_i$ gives
\[
 [F_j(U)]_{\{v\}}=\sum_i [D_{i,j}(U)\,stamp_i]_{\{v\}},
\]
which is the quoted singleton coefficient. The series converges absolutely:
the door-volume estimate bounds the hinge coefficients by the menu's
weighted load, and the menu offer has a summable fee-load ledger.

Group activity polynomials by the box label $p=\operatorname{barcode}(j)$:
\[
 K_p(U)=\sum_{j:\,\operatorname{barcode}(j)=p}F_j(U).
\]
Each fiber is countable. The paid coefficient majorants and their rooted
hit ledger imply absolute convergence on each fiber and preserve its price
when the terms are summed. The quiet coefficient of $K_p$ is the activity
profile on $p$: at the empty marker set the interpolated code is its vacuum
code, so its score reconstructs that value. Thus the profile term is
retained, while the singleton calculation above concerns only nonquiet
coefficients. Higher coefficients of $K_p$ contain all compatible products
and are not sums of singleton responses.

Write $\mathcal L_t(A,U)$ for the combined score of the full fragrance
response; its evaluation satisfies
$\sum_{A\subseteq F}t^{|A|}\mathcal L_t(A,U)=\sum_i\tau_i(U,t,F)$.
For any such score family $q_t(A)$, define its limiting marker polynomial by
\[
 \mathsf{Pass}(q)=\sum_A\left(\lim_{t\downarrow0}q_t(A)\right)
                           \prod_{i\in A}z_i.
\]
For a support mode $s$, define $\mathsf{Field}((K_p),s)$ by the finite sum
\[
\mathsf{Field}((K_p),s)=\sum_{\mathcal C:\,\bigsqcup_{p\in\mathcal C}Q_p=s}
                                  \prod_{p\in\mathcal C}K_p,
\]
where a box is a pair $p=(P_p,Q_p)$ consisting of its full support $P_p$
and assigned flag $Q_p\subseteq P_p$. In every family in the sum the full
supports $P_p$ are pairwise disjoint and the flags have union $s$; a
nonzero box term has nonempty assigned flag. This is the support pattern
condition: disjoint flags alone would permit two interacting polymers to
overlap and would not reproduce the gas. Grouping by barcode preserves the
condition. Indeed, two nonzero compatible activity terms cannot have the same
barcode, because that barcode has a nonempty full support. Thus each
barcode occurs at most once in a compatible family; its full support is
unchanged by grouping, and its assigned flag remains a subset of that
support. Expanding the products after this bijective grouping gives exactly
the original full-support-disjoint families and their flag union. Absolute
summability justifies the countable regrouping. Because the full logarithm score has zero quiet coefficient, its marker
polynomial is nilpotent on the finite marker algebra. Its truncated
exponential is therefore a finite polynomial. The full pushed expression is
\begin{equation}
 \operatorname{push}(\operatorname{Expense})(U)
 =e^{E(U)}\,\exp\!\bigl(\mathsf{Pass}(\mathcal L(U))\bigr)
  *\mathsf{Field}\bigl((K_p(U))_p,\operatorname{flags}(U)\bigr).
 \label{eq:ex-lattice-canopy-push}
\end{equation}
For a finite ticket truncation, expand the purchase as its finite ticket
sum and push each ticket term using linearity of the kernel integral; the
lawful input makes each term integrable. The score-to-marker map preserves
products because repeated switches vanish. Its exponential is the finite
nilpotent exponential, and it preserves the countable fragrance sum by the
summable majorants. Expanding the field pattern then gives
\eqref{eq:ex-lattice-canopy-push}; the absolute price bounds pass to the
countable limit. The quiet profile, singleton coefficient, support, shell,
measurability, and gauge conditions are exactly the receipts needed for the
per-box response. The hinge identity \eqref{eq:ex-lattice-menu-jet-input}
and the selected geometric price receipts are upstream inputs; the next
section proves the hinge identity itself.

\subsection{What these estimates produce}

The two row estimates have separate roles. The fragrance estimate gives a
summable family of paid logarithmic terms, with rooted budget
$2P_0/(1-\rho)$ at exponent $r_4$. The field estimate gives a paid
activity response at the straw budget \eqref{eq:ex-lattice-volant-straw}.
Together with the selected crest locality and the coefficientwise cash
identity \eqref{eq:ex-lattice-multi-closed}, they supply the paid-response
inputs to the per-box assembly in
Theorem~\ref{thm:ex-lattice-volant-supplier}. All numerical dependence on
$L,\beta,\mu$, the budget radius, and the menu profile remains visible in
the displayed prices.
